%======================================================================
\chapter{Security Analysis and Performance in the Base Paper}
\label{ch:basesec}
%======================================================================

Having set out what the protocol \emph{does}, this chapter covers the paper's argument that
it is \emph{secure}, and its accounting of what it costs. Three separate arguments are made:
a formal proof in BAN logic, a set of informal propositions, and machine verification with
Scyther.

\section{Threat model}

The paper assumes an adversary capable of:

\begin{itemize}
\item \textbf{Passive eavesdropping} --- recording everything on the channel.
\item \textbf{Active impersonation} --- claiming to be a legitimate node.
\item \textbf{Replay} --- retransmitting captured messages.
\item \textbf{Delay injection} --- deliberately slowing packets to push a legitimate
      exchange outside its timestamp window, causing desynchronisation.
\item \textbf{Physical tampering} --- compromising a subset of nodes and extracting their
      stored credentials.
\end{itemize}

The four defences claimed against these are: unique non-reusable nonces with adaptive
timestamp validation; asymmetric encryption with secure key derivation; tamper-resistant
storage with hierarchical credential management; and revocation combined with dynamic
fallback paths.

\section{Formal analysis with BAN logic}

BAN logic \cite{ban} is a formal system for reasoning about \emph{belief} in authentication
protocols. Rather than modelling bits, it models what each participant is entitled to
conclude after seeing a message. It is a proof technique, not a tool: you write down what
each party initially believes, apply inference rules to each message, and check whether the
desired beliefs follow.

\subsection{Notation}

For entities $\varepsilon, \phi$ and messages $M_\varepsilon, M_\phi$:

\begin{table}[H]
\centering
\small
\caption{BAN logic notation used in the base paper's proof.}
\label{tab:ban-not}
\begin{tabularx}{\textwidth}{@{}L{30mm} X@{}}
\toprule
\textbf{Notation} & \textbf{Reading}\\
\midrule
$\prec M_\varepsilon \succ$ & Message $M_\varepsilon$ is \emph{fresh}\\
$\varepsilon \Lleftarrow M_\varepsilon$ & $\varepsilon$ \emph{sees} the statement\\
$\varepsilon \Rightarrow M_\varepsilon$ & $\varepsilon$ has \emph{jurisdiction} over the statement\\
$\varepsilon \gg M_\varepsilon$ & $\varepsilon$ \emph{believes} the statement\\
$\varepsilon \sim M_\varepsilon$ & $\varepsilon$ \emph{once said} the statement\\
$\{M_\varepsilon\}_{\PU_\varepsilon}$ & Encrypted under $\varepsilon$'s public key\\
$\phi \xLeftrightarrow{\PU_\varepsilon} \varepsilon$ & The two communicate securely using $\PU_\varepsilon$\\
\bottomrule
\end{tabularx}
\end{table}

\subsection{Inference rules}

\begin{table}[H]
\centering
\small
\caption{The five BAN rules applied in the proof.}
\label{tab:ban-rules}
\begin{tabularx}{\textwidth}{@{}L{10mm} L{40mm} X@{}}
\toprule
& \textbf{Rule} & \textbf{Plain reading}\\
\midrule
R1 & Message-meaning &
  If I believe we share a key and I see a message encrypted under it, then I believe you
  sent that message.\\
R2 & Nonce-verification &
  If I believe the message is fresh and that you sent it, then I believe you currently
  believe it (not that you said it long ago).\\
R3 & Jurisdiction &
  If I believe you are authoritative over a statement and that you believe it, then I
  believe it.\\
R4 & Freshness &
  If any part of a message is fresh, the whole message is fresh.\\
R5 & Belief &
  If I believe you believe a compound message, I believe you believe each part.\\
\bottomrule
\end{tabularx}
\end{table}

\subsection{Goals, assumptions and the derivation}

Eight goals are stated, one per direction of each of the four hops:
\begin{align}
\Omega_1&: \mathrm{UWS} \gg (\mathrm{UWS} \xLeftrightarrow{\PU_{SUB}} \mathrm{SUB}) &
\Omega_2&: \mathrm{SUB} \gg (\mathrm{UWS} \xLeftrightarrow{\PU_{UWS}} \mathrm{SUB}) \notag\\
\Omega_3&: \mathrm{SUB} \gg (\mathrm{SUB} \xLeftrightarrow{\PU_{B}} \mathrm{B}) &
\Omega_4&: \mathrm{B} \gg (\mathrm{SUB} \xLeftrightarrow{\PU_{SUB}} \mathrm{B}) \notag\\
\Omega_5&: \mathrm{B} \gg (\mathrm{B} \xLeftrightarrow{\PU_{SAT}} \mathrm{SAT}) &
\Omega_6&: \mathrm{SAT} \gg (\mathrm{B} \xLeftrightarrow{\PU_{B}} \mathrm{SAT}) \notag\\
\Omega_7&: \mathrm{SAT} \gg (\mathrm{SAT} \xLeftrightarrow{\PU_{BS}} \mathrm{BS}) &
\Omega_8&: \mathrm{BS} \gg (\mathrm{SAT} \xLeftrightarrow{\PU_{SAT}} \mathrm{BS}) \notag
\end{align}

The initial assumptions are that every entity believes its own nonces and timestamps are
fresh ($A_1, A_2$), that each entity has jurisdiction over its own identity ($A_3$), and
that public keys are correctly distributed ($A_4$).

The derivation then runs for thirty steps. Figure~\ref{fig:ban} shows the first hop, which
is representative --- the other three hops repeat the identical five-step shape.

\begin{figure}[H]
\centering
\resizebox{\textwidth}{!}{%
\begin{tikzpicture}[x=1cm,y=1cm,font=\footnotesize]
\node[box,anchor=west,align=left,minimum width=68mm,fill=mist] (s1) at (0,0)
 {\textbf{s1} \ \ $\mathrm{SUB} \Lleftarrow \langle \ID_{UWS}\cat\Nonce_{UWS}\cat\TS_1\rangle_{\PU_{SUB}}$\\
  \textit{\scriptsize SUB sees the encrypted challenge \textemdash\ direct from $M_1$}};
\node[box,anchor=west,align=left,minimum width=68mm,fill=mist] (s2) at (0,-1.5)
 {\textbf{s2} \ \ $\mathrm{SUB} \gg (\mathrm{UWS} \sim \langle \ID_{UWS}\cat\Nonce_{UWS}\rangle)$\\
  \textit{\scriptsize by R1 + A4: only UWS could have produced it}};
\node[box,anchor=west,align=left,minimum width=68mm,fill=mist] (s3) at (0,-3.0)
 {\textbf{s3} \ \ $\mathrm{SUB} \gg \prec \ID_{UWS}\cat\Nonce_{UWS} \succ$\\
  \textit{\scriptsize by R4 + A1: the nonce is fresh, so the whole message is}};
\node[box,anchor=west,align=left,minimum width=68mm,fill=mist] (s4) at (0,-4.5)
 {\textbf{s4} \ \ $\mathrm{SUB} \gg \mathrm{UWS} \gg \langle \ID_{UWS}\cat\Nonce_{UWS}\rangle$\\
  \textit{\scriptsize by R2 on s2 and s3: UWS believes it \emph{now}, not historically}};
\node[box,anchor=west,align=left,minimum width=68mm,fill=palekelp,draw=kelp] (s5) at (0,-6.0)
 {\textbf{s5} \ \ $\mathrm{SUB} \gg (\mathrm{UWS} \xLeftrightarrow{\PU_{UWS}} \mathrm{SUB})$
  \ \ \textcolor{kelp}{\textbf{GOAL $\Omega_2$}}\\
  \textit{\scriptsize by R3 + A3: jurisdiction over its own identity closes the argument}};
\foreach \a/\b in {s1/s2,s2/s3,s3/s4,s4/s5} \draw[flow] (\a.south) -- (\b.north);
% right column: the mirror direction
\node[box,anchor=west,align=left,minimum width=68mm,fill=mist] (t1) at (8.2,-1.5)
 {\textbf{s6} \ \ $\mathrm{UWS} \Lleftarrow \langle \ID_{SUB}\cat\Nonce_{UWS}\cat\TS_2\rangle_{\PU_{UWS}}$\\
  \textit{\scriptsize UWS sees the acknowledgement \textemdash\ from $M_2$}};
\node[box,anchor=west,align=left,minimum width=68mm,fill=mist] (t2) at (8.2,-3.0)
 {\textbf{s7} \ \ $\mathrm{UWS} \gg (\mathrm{SUB} \sim \langle \ID_{SUB}\cat\Nonce_{UWS}\rangle)$\\
  \textit{\scriptsize by R1 + A4}};
\node[box,anchor=west,align=left,minimum width=68mm,fill=mist] (t3) at (8.2,-4.5)
 {\textbf{s8} \ \ $\mathrm{UWS} \gg \mathrm{SUB} \gg \langle \ID_{SUB}\cat\Nonce_{UWS}\rangle$\\
  \textit{\scriptsize by R2 \textemdash\ the echoed nonce is what proves liveness}};
\node[box,anchor=west,align=left,minimum width=68mm,fill=palekelp,draw=kelp] (t4) at (8.2,-6.0)
 {\textbf{s9} \ \ $\mathrm{UWS} \gg (\mathrm{UWS} \xLeftrightarrow{\PU_{SUB}} \mathrm{SUB})$
  \ \ \textcolor{kelp}{\textbf{GOAL $\Omega_1$}}\\
  \textit{\scriptsize by R3 + A3}};
\foreach \a/\b in {t1/t2,t2/t3,t3/t4} \draw[flow] (\a.south) -- (\b.north);
\node[font=\scriptsize\bfseries\sffamily,text=navy,anchor=west] at (0,0.95)
  {DIRECTION 1: SUB authenticates UWS \ (from $M_1$)};
\node[font=\scriptsize\bfseries\sffamily,text=navy,anchor=west] at (8.2,-0.55)
  {DIRECTION 2: UWS authenticates SUB \ (from $M_2$)};
\node[box,anchor=west,fill=paleamber,draw=amber,align=left,minimum width=150mm,font=\scriptsize] at (0,-7.4)
 {Both goals reached $\Rightarrow$ \textbf{mutual} authentication on hop 1. Steps s10--s30
  repeat this identical shape for SUB$\leftrightarrow$B ($\Omega_3,\Omega_4$),
  B$\leftrightarrow$SAT ($\Omega_5,\Omega_6$) and SAT$\leftrightarrow$BS ($\Omega_7,\Omega_8$).};
\end{tikzpicture}}
\caption[BAN logic derivation for hop 1]{\textbf{The BAN derivation, hop 1.} The left column
is the challenge direction, the right the acknowledgement direction. Reaching both goals is
what makes the authentication \emph{mutual}. The critical step is s4/s8, where the freshness
of the nonce upgrades ``$\phi$ once said this'' into ``$\phi$ believes this now'' ---
without freshness, a replayed message would satisfy the earlier steps.}
\label{fig:ban}
\end{figure}

\section{Informal security propositions}
\label{sec:informal}

The paper supplements the BAN proof with five propositions.

\begin{table}[H]
\centering
\caption{The paper's five security propositions and the mechanism behind each.}
\label{tab:props}
\small
\begin{tabularx}{\textwidth}{@{}L{8mm} L{34mm} X@{}}
\toprule
& \textbf{Claim} & \textbf{Mechanism}\\
\midrule
P1 & No identity spoofing &
  Each identity is bound to a key pair and validated by signature at registration;
  $\mathrm{Verify}(\mathrm{Sig}(\ID_i,\mathrm{SK}_i),\mathrm{PK}_i)$ must return true.
  Forging an identity means forging a signature.\\
\addlinespace[2pt]
P2 & Replay resistance &
  The receiver checks $N_i \notin \mathrm{UsedNonces}$ \emph{and}
  $|\mathrm{now}-\TS_i| \le \Delta T$. The paper further proposes an \emph{adaptive}
  $\Delta T$: each node keeps a rolling average of recent inter-node delays and adjusts the
  window, so legitimate variation is tolerated while stale injection is still caught.\\
\addlinespace[2pt]
P3 & No man-in-the-middle &
  Messages are encrypted under the recipient's public key. Only the holder of the matching
  private key can decrypt or meaningfully alter them.\\
\addlinespace[2pt]
P4 & Secure key establishment &
  $K_{ij} = \mathrm{DH}(\mathrm{PK}_i,\mathrm{SK}_j) = \mathrm{DH}(\mathrm{PK}_j,\mathrm{SK}_i)$.
  Only the two legitimate parties can compute it. Bounded $\Delta T$ per hop additionally
  causes a delay-injection attempt to abort the session.\\
\addlinespace[2pt]
P5 & Resists stolen-verifier and ephemeral-secret-leakage attacks &
  Long-term private keys authenticate session establishment only; session keys are
  ephemeral, derived from fresh nonces and timestamps, and never stored. Compromising a node
  therefore yields neither past nor future session keys --- this is forward secrecy in
  effect. The registration set is periodically checked against the BS's revocation list.\\
\bottomrule
\end{tabularx}
\end{table}

\section{Machine verification with Scyther}

Scyther \cite{scyther} is an automated model checker for security protocols developed by
Cas Cremers. You describe the protocol in SPDL (Security Protocol Description Language),
state the properties you want, and the tool searches exhaustively for a trace in which a
Dolev--Yao adversary violates one. If no such trace exists within the bound, it reports a
proof of correctness; if one does, it draws the attack.

The four properties checked are:

\begin{table}[H]
\centering
\caption{Scyther claim types and what each rules out.}
\label{tab:claims}
\small
\begin{tabularx}{\textwidth}{@{}L{24mm} X L{42mm}@{}}
\toprule
\textbf{Claim} & \textbf{Guarantee} & \textbf{Attack it rules out}\\
\midrule
\texttt{Secret} & The named value is never learned by the adversary:
  $\forall \mathcal{A},\ \Pr[\mathcal{A}(S)=S] \le \theta$ &
  Eavesdropping, key recovery\\
\texttt{Alive} & The claimed peer actually executed the protocol &
  Talking to a ghost\\
\texttt{Weakagree} & If $\varepsilon_A$ believes it ran with $\varepsilon_B$, then
  $\varepsilon_B$ really did run with $\varepsilon_A$:
  $P(\varepsilon_A,\varepsilon_B) \Rightarrow \exists P(\varepsilon_B,\varepsilon_A)$ &
  Man-in-the-middle, impersonation\\
\texttt{Niagree} / \texttt{Nisynch} & Both parties agree on all exchanged values, in a
  causally ordered sequence: a message received at $t_r$ was sent at $t_s \le t_r$ &
  Replay, message reordering, value substitution\\
\bottomrule
\end{tabularx}
\end{table}

The paper reports that all entities passed, configured for five roles with up to five
concurrent sessions each, across more than sixty test cases, with no attacks found. Its
Fig.~6 shows the tool output. The paper does not publish its SPDL source; our own
reconstruction, and the restructuring it needed, is the subject of Chapter~\ref{ch:scyther}.

\section{Performance in the base paper}

\subsection{Experimental setup}

The paper's evaluation uses a custom Python 3.10 simulator on Ubuntu 22.04, Intel Core
i7-11800H, 32\,GB RAM. The modelled network is a hierarchical tree of 10 to 200 nodes over 2
to 5 depth levels, with acoustic links at a nominal 10\,kbps, inter-node distances of 100 to
1000\,m and sound speed 1500\,m/s. Message fields are standardised at 64-bit identities,
64-bit nonces, 8-bit timestamps, 160-bit hashes and 512-bit ECC-point or AES-key
representations.

\subsection{Communication cost}

\begin{table}[H]
\centering
\caption[Communication cost comparison]{Communication cost per full authentication cycle,
in bits (base paper, Table~IV). ``Ref [21]--[25]'' are the scheme identifiers used in the
base paper's own comparison, and are \emph{not} the reference numbers of this report.}
\label{tab:comm}
\small
\begin{tabular}{@{}l r r r r r l@{}}
\toprule
\textbf{Protocol} & \textbf{Entity 1} & \textbf{Entity 2} & \textbf{Entity 3} &
\textbf{Entity 4} & \textbf{Total} & \textbf{vs proposed}\\
\midrule
Ref [21] & 896  & 768 & 512 & 832 & 3008 & $+42\,\%$\\
Ref [22] & 1024 & 832 & 512 & 832 & 3200 & $+52\,\%$\\
Ref [23] & 960  & 800 & 544 & 832 & 3136 & $+49\,\%$\\
Ref [24] & 912  & 784 & 512 & 832 & 3040 & $+44\,\%$\\
Ref [25] & 1008 & 848 & 528 & 832 & 3216 & $+52\,\%$\\
\midrule
\textbf{Proposed} & \textbf{640} & \textbf{512} & \textbf{320} & \textbf{640} &
\textbf{2112} & \textbf{baseline}\\
\bottomrule
\end{tabular}
\end{table}

The percentages are relative to the proposed scheme as denominator; for example
$(3008-2112)/2112 = 42\,\%$.

\subsection{Computational cost}

\begin{table}[H]
\centering
\caption[Computational cost comparison]{Computational cost per entity in milliseconds
(base paper, Table~V). $T_{exp}$ exponentiation, $T_{ecm}$ elliptic-curve multiplication,
$T_h$ hash, $T_{pair}$ bilinear pairing.}
\label{tab:comp}
\small
\begin{tabularx}{\textwidth}{@{}l r r X@{}}
\toprule
\textbf{Protocol} & \textbf{UWS / BS} & \textbf{SUB / SAT} & \textbf{Operations}\\
\midrule
Ref [21] & 0.536  & 0.800  & 2 exp $+$ 6 hash\\
Ref [22] & 19.70  & 25.000 & 3 ECM $+$ 4 ECM $+$ pairings\\
Ref [23] & 75.88  & 90.000 & 3 ECM $+$ 2 pairings\\
Ref [24] & 2.352  & 2.900  & 1 ECM $+$ 4 hash\\
Ref [25] & 1.245  & 1.600  & 1 ECM $+$ 3 hash\\
\midrule
\textbf{Proposed} & \textbf{0.400} & \textbf{0.500} &
\textbf{1 ECC $+$ 2 hash $+$ 4 AES}\\
\bottomrule
\end{tabularx}
\end{table}

The saving comes from eliminating bilinear pairings and multiple exponentiations --- the two
operations that dominate cost in the older schemes --- and replacing them with one optimised
ECC operation plus cheap symmetric work. Against Ref~[23] this is a $189\times$ improvement.

\subsection{Energy per authentication cycle}

The paper's energy accounting uses published ARM Cortex-M4 / Cortex-A53 benchmarks:
one ECC operation $\approx 24\,\mu$J, one hash $\approx 6\,\mu$J, one AES-128
encryption or decryption $\approx 3.2\,\mu$J. A single entity performs one ECC operation,
two hashes, and two encryptions plus two decryptions:

\begin{equation}
E_{cycle} = \underbrace{24}_{\text{ECC}}
          + \underbrace{2 \times 6}_{\text{hash}}
          + \underbrace{(2+2) \times 3.2}_{\text{AES}}
          = 24 + 12 + 12.8 = 48.8\ \mu\mathrm{J}
\label{eq:energy}
\end{equation}

\begin{table}[H]
\centering
\caption{Energy per authentication cycle against prior schemes.}
\label{tab:energy}
\small
\begin{tabular}{@{}l r l@{}}
\toprule
\textbf{Protocol} & \textbf{Energy} & \textbf{Relative}\\
\midrule
Ref [22] & $\approx 2300\ \mu$J & $47\times$ more\\
Ref [23] & $\approx 8900\ \mu$J & $182\times$ more\\
Ref [24] & $\approx 280\ \mu$J  & $5.7\times$ more\\
\midrule
\textbf{Proposed} & $\bm{48.8\ \mu}$\textbf{J} & \textbf{baseline}\\
\bottomrule
\end{tabular}
\end{table}

At 48.8\,\uJ{} --- under 0.05\,mJ --- a node with a 1000\,mAh cell at 3.3\,V
(3.3\,J $=3.3\times10^{6}$\,\uJ) could in principle perform about 67\,000 authentication
cycles on cryptography alone. Chapter~\ref{ch:results} shows what happens once transmission
and reception energy are added, which is where the real budget goes.
